\documentclass[12pt]{article}
\usepackage{amsmath,amssymb}

\begin{document}
\section*{{2023 AMC 12A Problems/Problem 21}}
Source: AoPS Wiki

\subsection*{Solution}
To find the total amount of vertices we first find the amount of edges, and that is $\frac{20 \times 3}{2}$ . Next, to find the amount of vertices we can use Euler's characteristic, $V - E + F = 2$ , and therefore the amount of vertices is $12$ So there are $\binom{12}{3} = 220$ ways to choose 3 distinct points. Now, the furthest distance we can get from one point to another point in an icosahedron is 3. Which gives us a range of $1 \leq d(Q, R), d(R, S) \leq 3$ With some case work, we get two cases: Case 1: $d(Q, R) = 3; d(R, S) = 1, 2$ Since we have only one way to choose Q, that is, the opposite point from R, we have one option for Q and any of the other points could work for S. Then, we get $12 \times 1 \times 10 = 120$ (ways to choose R × ways to choose Q × ways to choose S) Case 2: $d(Q, R) = 2; d(R, S) = 1$ We can visualize the icosahedron as 4 rows, first row with 1 vertex, second row with 5 vertices, third row with 5 vertices and fourth row with 1 vertex. We set R as the one vertex on the first row, and we have 12 options for R. Then, Q can be any of the 5 points on the third row and finally S can be one of the 5 points on the second row. Therefore, we have $12 \times 5 \times 5 = 300$ (ways to choose R × ways to choose Q × ways to choose S) In total, we have $420$ ways, but we must divide by $3!$ to account for permutations, giving us $70/220 = \boxed{\textbf{(A) } 7/22}.$ ~lptoggled, edited by grogg007 , ESAOPS, weimi23, and trevian1

\end{document}
